\chapter{Conclusion et perspectives}
\label{ch:ccl}

\section{Bilan}
POSTIE démontre la faisabilité d'un copilote PO open-source
articulant RAG multi-modèles et priorisation algorithmique
multi-critères. Le projet a permis de livrer un produit fonctionnel
en six sprints, de formaliser une démarche Agile complète et
d'identifier les bénéfices comme les limites concrètes de l'IA
générative appliquée au product management. La mise en production
effective (frontend Vercel, backend Render) prouve la maturité
suffisante du prototype pour une démonstration en condition
réelle.

\section{Perspectives}
\begin{itemize}
  \item Migration vers une base \textbf{Postgres} persistante
        (Neon, free tier) pour le backlog et l'historique
        conversationnel.
  \item Intégration native \textbf{Jira / Linear / Azure DevOps}
        via OAuth pour synchroniser bidirectionnellement le backlog
        avec les outils existants de l'équipe.
  \item Évaluation longitudinale sur une équipe produit réelle
        (mesure de vélocité, qualité des stories, satisfaction).
  \item Entraînement d'un classifieur dédié pour remplacer le
        routage par heuristique.
  \item Étude d'un mode auto-hébergé (Mistral / Llama 3) pour les
        contextes contraints en confidentialité (santé, défense,
        finance).
  \item Mode multi-utilisateurs avec gestion de rôles
        (PO, Scrum Master, Stakeholder).
\end{itemize}

\section{Apports personnels}
Ce projet m'a permis de consolider mes compétences en architecture
logicielle full-stack, en orchestration LLM moderne (LangChain,
multi-modèles, RAG), en pilotage Agile autonome d'un produit du
\gls{mvp} jusqu'à la mise en production, et en communication
technique vulgarisée envers les parties prenantes non techniques.

Au-delà des compétences techniques, POSTIE m'a confronté aux
arbitrages réels du métier de PO~: prioriser sous incertitude,
défendre des choix devant un manager, livrer en respectant un
échéancier serré. Cette \emph{double posture} --- développer
l'outil tout en étant son utilisateur cible --- a constitué l'un
des apports pédagogiques majeurs de ce mémoire.
